\section{Conclusion}

Answer-visible process verification should not be trusted on ordinary
discrimination alone. Across executable matched-quartet benchmarks,
answer-matched counterfactual audits reveal a stable gap between ordinary
separability and genuinely local process discrimination, and mechanism analyses
support a picture in which lower local \amcd{} accompanies the downstream
utility cost of answer sensitivity.

The deployment takeaway is therefore simple: in the current setting,
the safest intervention is to remove direct answer access. On
\textsc{Craft-Deploy}, masked reranking clearly outperforms visible reranking;
on \textsc{Craft-Clean}, masking still improves process-faithfulness even when
utility no longer separates. The most durable contribution is thus an
audit-and-deployment recipe for process verifiers, not a claim that a single
visible repair term already solves answer leakage. The external-source
\textsc{ProcessBench} package reinforces the same conservative read on real
step-level traces rather than overturning it. The broader lesson is that
process verifiers should be validated under answer-matched counterfactual
audits before they are trusted in reranking or compute-allocation loops.
